% !PRINTED 28
\weberpage{28}

%<<CONT
wenn es zu $P_\nu$ mehrere Theiler von der Beschaffenheit wie $Q_\nu$
giebt, dass einer von möglichst niedrigem Grade für $Q_\nu$ genommen
 ist.

Wenn wir dann nachweisen, wie man aus $Q_\nu$ einen commutativen
 Normaltheiler $Q_{\nu-1}$ von $P_{\nu-1}$ ableiten kann, der zugleich
 $Q_\nu$ enthält, so können wir dies Verfahren fortsetzen, bis
wir zu einem commutativen Normaltheiler $Q$ von $P$ selbst gelangt
 sind, wie ihn unser Satz verlangt.

Wählen wir irgend ein Element $\gamma$ in der Gruppe $P_{\nu-1}$,
welches nicht in $P_\nu$ enthalten ist, so ist, da $P_\nu$ ein Normaltheiler
von $P_{\nu-1}$ ist,

\weq{3}{\gamma P_\nu = P_\nu \gamma,}

und wenn daher $h$ der niedrigste Exponent ist, für den $\gamma^h$ in
$P_\nu$ enthalten ist, so ist $h$ grösser als 1 und

\weq{4}{\gamma^h P_\nu = P_\nu,}

und die in dem System der Nebengruppen

\weq{5}{P_\nu,\ \gamma P_\nu,\ \gamma^2 P_\nu \ldots \gamma^{h-1} P_\nu}

enthaltenen Elemente, die alle von einander verschieden sind,
bilden eine Gruppe, deren Grad, wenn der Grad von $P_\nu$ mit $p_\nu$
bezeichnet wird, gleich $h p_\nu$ ist. Die Gruppe (5) ist aber auch
ein Theiler von $P_{\nu-1}$ und folglich muss $h p_\nu$ ein Theiler von
$p_{\nu-1} = j_\nu p_\nu$ sein, und $h$ ist ein Theiler von $j_\nu$. Da aber $j_\nu$
eine Primzahl ist, so muss $h = j_\nu$ sein, und durch (5) ist die
ganze Gruppe $P_{\nu-1}$ erschöpft:

\weq{6}{P_{\nu-1} = P_\nu + \gamma P_\nu + \gamma^2 P_\nu + \cdots \gamma^{j_\nu - 1} P_\nu.}

Wir betrachten nun das System der transformirten Gruppen

\weq{7}{\gamma^{-1} Q_\nu \gamma,}

worin $\gamma$ alle Elemente von $P_{\nu-1}$ durchläuft.

Diese Gruppen sind alle mit einander isomorph und sind
also ebenso wie $Q_\nu$ commutativ. Sie sind alle in $P_\nu$ enthalten,
und sind, da $\gamma^{-1} P_\nu \gamma = P_\nu$ ist, nach dem Schlusssatze von \S.\,3
Normaltheiler von $P_\nu$.

Wenn alle Gruppen (7) mit einander identisch sind, so ist
$Q_\nu$ Normaltheiler von $P_{\nu-1}$, und wir erreichen unseren Zweck
der Bestimmung von $Q_{\nu-1}$ schon dadurch, dass wir $Q_\nu$ selbst,
oder, falls noch ein commutativer Normaltheiler niedrigeren
Grades von $P_{\nu-1}$ vorhanden sein sollte, diesen letzteren, für $Q_{\nu-1}$
nehmen.
